\section{Introduction}\label{sec:introduction}

How large can a convex body's symplectic capacity be if its volume is fixed?
Viterbo's conjecture proposed that a ball gives the largest value
\cite{Viterbo2000}. For the Ekeland--Hofer--Zehnder capacity in dimension
four, the proposed inequality was
\[
 \operatorname{sys}(K):=\frac{c_{\rm EHZ}(K)^2}{2\operatorname{vol}_4(K)}\leq1.
\]
The quantity on the left is called the systolic ratio. Capacity scales by
the square of a dilation and four-dimensional volume by its fourth power,
so this ratio compares shape independently of overall size. Translations
and linear symplectic maps also preserve it.

The inequality is false. Haim--Kislev and Ostrover found a counterexample
formed from two regular pentagons in complementary Lagrangian planes
\cite{HaimKislevOstrover2024}. Writing the coordinates as $(q,p)\in
\mathbb R^2\oplus\mathbb R^2$, one representative is
\[
 K_{\rm HKO}=P_5\times R(-\pi/2)P_5,\qquad
 \operatorname{sys}(K_{\rm HKO})=\frac{3+\sqrt5}{5}>1.
\]
Here $P_5$ is a centred regular pentagon with a vertex on the positive
horizontal axis, and $R(\theta)$ is planar rotation. A product placed in
the $q$- and $p$-planes is called a Lagrangian product. The relative
position of its factors affects capacity even though it does not affect
volume.

The counterexample leaves several geometric questions. Is the HKO body a
local maximum of the systolic ratio among convex bodies? Are there other
families with large ratios? Can the
computations reveal geometric relationships involving capacity? This thesis
investigates these questions through local analysis, finite algorithms and
data science. This is the sense in which we continue to probe Viterbo's
conjecture: although its proposed bound is false, the geometric relationship
between symplectic capacity and volume remains to be understood.

The mathematical results identify two ways of changing the counterexample
that cannot improve its ratio: sufficiently small ten-facet perturbations
and arbitrary independent linear deformations of its regular-pentagon
factors. These are different restrictions, one local in the full ten-facet
space and one global within a product family. Finite capacity formulas make
both questions accessible to proof. The empirical investigation asks whether
geometric measurements can explain capacity in broader populations and guide
the search for further examples.

\subsection{Capacities, dynamics and convex geometry}

Symplectic capacities measure constraints on embeddings that volume alone
cannot detect. Two basic examples are the Gromov width $c_G$, defined by
the largest symplectic ball that embeds into a domain, and the cylindrical
capacity $c_Z$, defined by the smallest symplectic cylinder into which the
domain embeds. With the usual normalization, every symplectic capacity lies
between these two~\cite{Edtmair2024}. The EHZ capacity has a different
description: on convex bodies it is the least action of a closed
characteristic. Thus comparing capacities connects an embedding question
to a question about Hamiltonian trajectories.

The conjecture that all normalized capacities agree on convex domains was
stronger than Viterbo's volume bound. Indeed, the Gromov width satisfies
that bound directly by volume preservation. The HKO example therefore
also separates the EHZ capacity from the Gromov width
\cite{HaimKislevOstrover2024}. Abbondandolo, Edtmair and Kang prove that EHZ and cylindrical capacity
coincide on every convex body in four dimensions
\cite[Corollary~1]{AbbondandoloEdtmairKang2024}. Thus the capacity computed
in this thesis also determines cylindrical capacity. Their proof builds on
Edtmair's earlier work, which characterizes
cylindrical capacity on smooth star-shaped domains in $\mathbb R^4$ with
dynamically convex boundary by the
least area of a disk-like global surface of section, and proves coincidence
of all normalized capacities for convex domains sufficiently $C^3$-close
to the ball~\cite{Edtmair2024}. A global surface of section is a surface
whose interior is transverse to the flow and is met in both forward and
backward time by every trajectory outside its boundary. These results illustrate how information about the
flow can determine an embedding obstruction.

There is also a direct connection to a central problem in convex geometry.
For a centrally symmetric convex body $L\subset\mathbb R^n$ and its polar
$L^\circ$, Artstein--Avidan, Karasev and Ostrover prove
$c_{\rm EHZ}(L\times L^\circ)=4$~\cite{AKO2014}. Viterbo's proposed
inequality on these particular products is consequently equivalent to the
lower bound $\operatorname{vol}(L)\operatorname{vol}(L^\circ)\geq4^n/n!$
in the symmetric Mahler conjecture. The general counterexample does not
invalidate this equivalence for the restricted family. Capacity--volume
questions therefore retain significance beyond the original universal bound.

Local maximality has a dynamical precedent as well. Abbondandolo and
Benedetti prove that Zoll contact forms, whose Reeb orbits are all closed
with the same minimal period, locally maximize the contact systolic ratio
\cite{AB2023}. Their results include the volume--capacity inequality near
the ball. The HKO body presents a different setting: it has corners and a
ratio larger than the ball's. Our local theorem studies how such a
nonsmooth example behaves under perturbations of its supporting facets.

Haim--Kislev studies the HKO body through hyperplane cuts. A convex body is
a \emph{cuts extremizer} if every cut into two pieces with nonempty interior
strictly decreases the EHZ capacity on each piece. Numerical experiments
suggest this property for HKO, while also suggesting that its capacity is
not always the sum of the capacities of the pieces
\cite[Definitions~1.1--1.2 and Example~1.12]{HaimKislev2025Zoll}.
For sufficiently shallow cuts in certain directions, including all normals
supported in one pentagon factor, Haim--Kislev proves this additivity identity
\cite[Proposition~1.13]{HaimKislev2025Zoll}.

These cuts give a complementary way to perturb the counterexample. Cutting
off a small cap typically adds a facet, whereas our theorem varies the ten
existing facets. The systolic ratio also compares capacity with volume: for
a retained piece $K'$, the inequality
$\operatorname{sys}(K')\leq\operatorname{sys}(K_{\rm HKO})$ is equivalent to
\[
 \frac{c_{\rm EHZ}(K')}{c_{\rm EHZ}(K_{\rm HKO})}
 \leq
 \left(\frac{\operatorname{vol}_4 K'}
 {\operatorname{vol}_4 K_{\rm HKO}}\right)^{1/2}.
\]
Thus the capacity decrease must be large enough to compensate for the lost
volume. The theorem below establishes this balance throughout a neighborhood
in the space of ten-facet polytopes.

\subsection{A local theorem around the counterexample}

Our main result partially answers the first question: every sufficiently
Hausdorff-close ten-facet polytope has ratio at most
that of $K_{\rm HKO}$. In a smaller neighborhood, equality occurs precisely
for translations, positive dilations and linear symplectic images of
$K_{\rm HKO}$; see Theorem~\ref{thm:hko-ten-facet-local-maximum}. Thus a
small change of shape that is not one of these symmetries lowers the ratio.
Allowing facets to appear or disappear is a different local problem, which
the theorem leaves open.

The proof uses upper bounds that agree with the systolic ratio at HKO.
For every nonzero direction transverse to the symmetries, at least one of
these bounds has negative derivative. A uniform estimate then gives strict
decrease away from HKO on a sufficiently small transverse slice.
The finite task is to construct enough
such bounds and check their derivatives. Feasible weights in Haim--Kislev's
capacity formula provide the bounds. An algebraic computation establishes the derivative
conditions; the proof explains why those conditions imply the geometric
local maximum.

\subsection{From boundary dynamics to finite calculations}

For a smooth convex body, the EHZ capacity is the least action of a closed
characteristic on its boundary. Polytopes have corners, so the relevant
curves are generalized characteristics. Their velocities are governed by
the normals of the facets they meet. Haim--Kislev's theorem makes the
capacity computable by a finite optimization over facet orders and
nonnegative weights~\cite{HK2017}.

We derive the finite formula from the dual action principle, expanding
Haim--Kislev's argument for a master's-level reader. We state the imported
variational results and explain the existence of a minimizing curve with a
finite list of distinct facet directions.

For Lagrangian products, we derive a further reduction from the bilinear
structure of the objective: some maximizing word uses at most three facets
from each factor. This gives a smaller finite search for the capacity,
without classifying every minimizing orbit.

Using the local passage geometry of Chaidez and Hutchings~\cite{CH2021},
we develop a flow-graph algorithm with a finite search justified by the
simple-minimizer theorem. Its implementation provides an independent
cross-check of the quadratic-program method on selected rational inputs
satisfying the flow algorithm's regularity hypotheses.

For regular pentagon products, the finite formula also gives the complete
relative-rotation profile. The systolic ratio of
$P_5\times R(\theta)P_5$ is
\[
 \frac{5+2\sqrt5}{10}\sec^2 d(\theta),
\]
where $d(\theta)$ is the distance to the nearest multiple of $\pi/5$.
A short analytic comparison proves the formula, whose maximum is the HKO
value. A separate matrix-containment argument extends the capacity formula
and sharp systolic bound to independent linear deformations of both regular
pentagons; see Section~\ref{sec:pentagon-affine}. This family is larger
than the rotation family but does not include arbitrary pentagons.

We also vary the placement of the factor planes themselves. For any two
full-dimensional planar convex bodies, a small rigid rotation away from the
symplectic-product arrangement increases capacity. The Lagrangian endpoint
need not be better: an explicit triangle product has smaller capacity there;
see Section~\ref{sec:product-rotation}.

\subsection{Looking for mathematical structure in data}

Symplectic ridge area gives a strong inverse association with the computed
systolic ratio in the sampled populations. Selection based on geometric
measurements also enriches fresh samples. Yet the overall rank association
is substantially weaker in subsets selected within each sampling group for
either low ridge scores or high computed ratios. Its success as a coarse
predictor leaves open how to improve a body already selected for a promising
score.

Affine shape variation of the factors contributes to the broad association.
Area-preserving linear maps that make each factor's vertex covariance
isotropic raise the computed ratio in 314 of 320 tested products and
substantially weaken the ridge--ratio correlation. This identifies a useful
shape deformation, although its effects cannot be attributed to covariance
imbalance alone: the maps change several geometric features together.

Exact family laws further constrain the explanation. An edge--width formula
describes the ridge measurement on Lagrangian products. Under relative
rotation of regular pentagon factors, it has an inverse-square relation to
the systolic ratio; under relative rotation of equilateral-triangle factors,
the two quantities have the same strict ordering; see
\eqref{eq:ds-pentagon-ridge-capacity} and~\eqref{eq:ds-triangle-profile}.
A general explanation must therefore account
for the shapes and relative positions of the factors.

Direct local search gives a separate positive result. In the comparison of
fixed implementations from the same 64 starting bodies, retaining previously
active capacity branches attains larger median best ratios than literal
gradient ascent under the same evaluation limits. This connects the
first-order branch model to practical improvement; the finite searches still
do not establish local maximality.

\subsection{How to read the thesis}

The foundations in Sections~\ref{sec:preliminaries}--\ref{sec:quadratic-program-algorithm-hk2019}
explain how nonsmooth boundary dynamics become a finite quadratic program.
Section~\ref{sec:flow-graph-algorithm-ch2021} presents the alternative flow-graph
method, and Section~\ref{sec:first-order-perturbations} develops first-order variation.

A reader primarily interested in the HKO theorem can start with its statement
and proof strategy in Section~\ref{sec:hko-local-maximum}, returning to the
foundations for the capacity formula and to
Section~\ref{sec:first-order-perturbations} for variation. The flow-graph and
data-science chapters are not prerequisites for that proof.
Section~\ref{sec:black-box-datascience} describes the empirical investigation.
Sections~\ref{sec:rotated-regular-polygons} and~\ref{sec:pentagon-affine} give
the rotation profile and its extension to independently deformed pentagons.
Section~\ref{sec:product-rotation} then rotates the full product out of
symplectic position, rather than rotating one factor within its plane.
The later sections document visualizations, numerical methods, code and data,
and reflect on the use of AI. The appendices contain computational detail
for checking or reproducing the results.
